\addtocontents{toc}{\protect\Needspace{3\baselineskip}}
\originalchapter{早期宇宙测验习题}
以下题面来自 8.286 Fall 2013 的正式测验, 不是复习题库. Quiz 1 第 3 题与 Quiz 3 第 2 题的官方答案仅回指作业答案, 未在该解答文件展开; 两题明确由 AI 独立补写. 其余六题依据官方解答核对并重新展开.

\src{8.286 Quiz 1, 2013年10月3日, Q2, (a)--(e), 原题 30 分; 官方解答标题误标 25 分.}
\begin{exercise}
平直宇宙 $a(t)=a_0\ee^{\chi t}$, 取 $\chi>0$. 两共动星系在 $t_1$ 的实际间距为 $\ell_p$, 一个在 $t_1$ 向另一个发光, 接收时刻 $t_2$.
(a) 求红移; (b) 仅以 $c,a_0,\chi,\ell_p,t_1$ 求 $t_2$; (c) 什么间距使光永不收到?
(d) 光在什么时候距两星系相等?
(e) 发射脉冲时长 $\Delta t$ 很短, 求收到的脉冲时长.
\end{exercise}
\begin{solution}
(a) $1+z=a(t_2)/a(t_1)=\ee^{\chi(t_2-t_1)}$.
(b) 共动间距 $\ell_c=\ell_p/a(t_1)$, 光传播给 $\ell_c=\int_{t_1}^{t_2}c\dd t/a(t)$. 因而 $\ell_p=(c/\chi)[1-\ee^{-\chi(t_2-t_1)}]$, 解得 $t_2=t_1-\chi^{-1}\ln(1-\chi\ell_p/c)$.
(c) 当 $\ell_p\ge c/\chi$ 时没有有限 $t_2$, 等号情形仅在无穷未来逼近.
(d) 同时的实际距离均含同一 $a(t)$, 中点因此也是共动中点. 积分路径改为 $\ell_c/2$, 得 $t_{\rm eq}=t_1-\chi^{-1}\ln[1-\chi\ell_p/(2c)]$. 中点能否到达的条件是 $\ell_p<2c/\chi$, 所以某些光永不到终点但可先过中点.
(e) 固定共动路径对端点微分给 $c\Delta t/a(t_1)=c\Delta t_o/a(t_2)$, 故 $\Delta t_o=(1+z)\Delta t$. 时长变长不是本地仪器走慢, 而是两次信号连接的宇宙时刻变化.
\end{solution}

\src{8.286 Quiz 1 Q3, (a)--(e), 25 分; 独立补写.}
\begin{exercise}
无限长均匀圆柱宇宙只沿 $x,y$ 膨胀, $z$ 方向静止. 初始密度 $\rho_i$, 初始速度 $\boldsymbol v_i=H_i\boldsymbol r$, 初始外半径有限. $\mu$ 为半径 $r$ 内每长度质量.
(a) Gauss 引力定律给 $\boldsymbol g=-A\mu\hat{\boldsymbol r}/r$, 求 $A$.
(b) 以初始半径 $r_i$ 写粒子 $\ddot r$.
(c) 证明 $u=r/r_i$ 与 $r_i$ 无关, 定义 $a(t)=u$.
(d) 求密度与尺度因子关系及 $\ddot a$.
(e) 找守恒量 $\mathcal E=\dot a^2/2+V(a)$, 判断永远膨胀还是再坍缩.
\end{exercise}
\begin{solution}
对单位长圆柱, $g_r2\pi r=-4\pi G\mu$, 故 $A=2G$. 未壳交叉时每层 $\mu=\pi\rho_ir_i^2$, $\ddot r=-2\pi G\rho_ir_i^2/r$. 代 $r=r_iu$ 得 $\ddot u=-2\pi G\rho_i/u$, 且对所有非中心层 $u(t_i)=1,\dot u(t_i)=H_i$. 正半轴的方程局部唯一, 相同初值因此给相同函数; 中心层由连续极限处理.

横截面积按 $a^2$ 变, $\rho=\rho_i/a^2$, 因而 $\ddot a=-2\pi G\rho a$. 乘 $\dot a$ 积分得 $V(a)=2\pi G\rho_i\ln a$、$\mathcal E=H_i^2/2$. 若初始膨胀, 转向点 $a_{\max}=\exp[H_i^2/(4\pi G\rho_i)]$ 有限; 对正密度和有限初始速度, 不存在一直向无穷半径逃逸的解. 对数势在大 $a$ 无界, 所以最终回缩; 这是无限长圆柱的性质, 不能套用到三维球宇宙.
\end{solution}

\src{8.286 Quiz 1 Q4, (a)--(b), 25 分.}
\begin{exercise}
今日宇宙时刻为 $t_0$, 星系当前实际距离 $\ell_p=a(t_0)\ell_c$. 视它为球, 发光时半径 $R_1$, 今日半径 $R_0$, 红移已知为 $z$, $\theta\ll1$.
(a) 不预设 $a(t)$ 的具体形式, 求张角.
(b) 发射总功率 $P$ 各向同性, 求今日地球收到的能量通量 $J$. 原题未写平直假设, 官方解答说明按当时课程的平直几何处理.
\end{exercise}
\begin{solution}
采用官方明确的平直假设. 星系发射时的共动半径是 $R_1/a_e$, 因此直径张角 $\theta=2R_1/(a_e\ell_c)=2R_1(1+z)/\ell_p$. 不应把 $R_0$ 代替 $R_1$, 引力束缚星系并不随 Hubble 流按统一比例膨胀.

今日波面面积 $4\pi\ell_p^2$, 每光子能量与到达率各降 $1+z$, 故 $J=P/[4\pi\ell_p^2(1+z)^2]$. 若空间弯曲, 应用横向距离 $D_M$ 替代简单 $\ell_p$, 题面原量不足以单独给此形式.
\end{solution}

\src{8.286 Quiz 2, 2013年11月7日, Q2, (a)--(d), 20 分.}
\begin{exercise}
平直宇宙由奇异物质主导, $\rho\propto a^{-5}$.
(a) 求尺度因子随时间的关系; (b) 求 $H(t)$; (c) 求实际粒子视界; (d) 求 $\rho(t)$.
\end{exercise}
\begin{solution}
从大爆炸起点 $t=0$ 计时, $\dot a/a\propto a^{-5/2}$, 故 $\dot a\propto a^{-3/2}$. 积分给 $a=Ct^{2/5}$, $H=2/(5t)$. 粒子视界 $D_p=a(t)c\int_0^t\dd t'/a(t')=5ct/3$. 由 Friedmann 式, $\rho=3H^2/(8\pi G)=3/(50\pi Gt^2)$. 能量守恒所对应状态参数为 $w=2/3$, 不等于通常辐射. 每项都适用此理想单组分模型, 不作为真实宇宙的历史.
\end{solution}

\src{8.286 Quiz 2 Q4, (a)--(c), 20 分.}
\begin{exercise}
假想物质能量密度随活塞体积 $V$ 按 $V^{-3/2}$ 下降, 初始为 $u_0$.
(a) 活塞拉大 $\Delta V$ 后, 求系统能量变化 $\Delta U$, 增加规定为正.
(b) 压强为 $p$, 求拉活塞的外界对系统做功 $\Delta W$.
(c) 用两者确定 $p$ 与 $u$ 的关系. 忽略二阶小量.
\end{exercise}
\begin{solution}
密度不是总能量. $U(V')=u_0(V/V')^{3/2}V'$, 在 $V'=V$ 求导给 $\Delta U=-(1/2)u_0\Delta V$. 外界对系统做功为 $\Delta W=-p\Delta V$. 没有额外热交换时两者相等, 故 $p=u/2$. 若把 $u\propto V^{-3/2}$ 直接当作 $U$ 的变化, 会错误得到 $3u/2$; 少掉的因子正是体积本身.
\end{solution}

\src{8.286 Quiz 3, 2013年12月5日, Q2, (a)--(c), 25 分; 独立补写.}
\begin{exercise}
缪子静止能 $106\,\mathrm{MeV}$, 电子 $0.511\,\mathrm{MeV}$, 两者均自旋 $1/2$, 缪子与反缪子电荷相反. 本题按理想阈值热辐射模型, 不考虑额外强子热贡献.
(a) 缪子与反缪子对 $u=(\pi^2/30)g(\kb T)^4/(\hbar c)^3$ 的 $g$ 贡献是多少?
(b) 当 $\kb T$ 刚高于 $106\,\mathrm{MeV}$ 时还有哪些轻粒子, 各贡献多少?
(c) 缪子消失时所有其他粒子仍保持共同热平衡, $aT$ 增加多少倍?
\end{exercise}
\begin{solution}
两电荷态各两自旋态, 费米权重 $7/8$, 所以 $g_\mu=4(7/8)=7/2$. 其余为光子两偏振 $g_\gamma=2$, 电子正电子 $g_e=7/2$, 三味中微子及反中微子共六个手征自由度 $g_\nu=6(7/8)=21/4$. 合计不含缪子 $g_{
\rm after}=43/4$, 含缪子 $g_{\rm before}=57/4$. 在全部共同温度且熵绝热守恒的模型中,
$[aT]_{\rm after}/[aT]_{\rm before}=(57/43)^{1/3}=1.0985$.
实际 $\kb T\sim106\,\mathrm{MeV}$ 并非严格的阶跃阈值, 接近质量阈值的粒子贡献连续变化, QCD 热史也更复杂. 原题的无中间质量叙述支持这一课堂自由粒子近似, 不把结果宣称为精确现代热史.
\end{solution}

\src{8.286 Quiz 3 Q3, (a)--(c), 25 分.}
\begin{exercise}
平直宇宙今日参数 $\Omega_{m0},\Omega_{r0},\Omega_{\Lambda0},H_0$, 三密度参数和为 1. 今日从原点发出球面光脉冲.
(a) 已知 $a(t)$ 时写永远能逼近的最大共动半径 $r_{\max}$.
(b) 换积分变量, 仅以今日参数和 $a_0$ 写此积分, 不要求求值.
(c) 对今天收到的光, 若源的共动位置恰在 $r_{\max}$, 写求源红移 $z_{\rm eh}$ 的显式积分方程.
\end{exercise}
\begin{solution}
零测地线给 $\dot r=c/a$. 最大半径为 $r_{\max}=c\int_{t_0}^\infty\dd t/a(t)$; 此积分需由未来模型保证收敛. 令 $x=a/a_0$, $H=H_0\sqrt{\Omega_{r0}x^{-4}+\Omega_{m0}x^{-3}+\Omega_{\Lambda0}}$, 则
\[
r_{\max}=\frac c{a_0H_0}\int_1^\infty
\frac{\dd x}{\sqrt{\Omega_{r0}+\Omega_{m0}x+\Omega_{\Lambda0}x^4}}.
\]
今天收到的光的过去传播距离也必须是这个共动数, 所以
\[
\int_0^{z_{\rm eh}}\frac{\dd z}{\sqrt{\Omega_{r0}(1+z)^4+\Omega_{m0}(1+z)^3+\Omega_{\Lambda0}}}
=\int_1^\infty\frac{\dd x}{\sqrt{\Omega_{r0}+\Omega_{m0}x+\Omega_{\Lambda0}x^4}}.
\]
「该源今天的位置在事件视界」不表示它的早期光永远不可见. 示例忽略辐射、取 $0.3,0.7,H_0=71$, 今日实际事件视界约 $4816\,\mathrm{Mpc}$, 对应今天收到的源红移约 $1.805$. 这只是给定未来真空不变模型的演示.
\end{solution}

\src{8.286 Quiz 3 Q4, (a)--(b), 15 分.}
\begin{exercise}
(a) 从 $H^2=8\pi G\rho/3-kc^2/a^2$ 推导 $(\Omega-1)/\Omega=A T^2/\rho$, 找常数 $A$.
(b) 对近乎平直、$\rho\propto a^{-5}$ 的假想宇宙, 求该偏离量随时间的变化. 温度作为光子温度, 假设 $aT=C$ 固定.
\end{exercise}
\begin{solution}
$\Omega=8\pi G\rho/(3H^2)$, 所以 $(\Omega-1)/\Omega=1-3H^2/(8\pi G\rho)=3kc^2/(8\pi G\rho a^2)$. 用 $aT=C$ 得 $A=3kc^2/(8\pi GC^2)$. 常数依尺度因子的规范与温度归一化而变, 整个偏离量不变. 近乎平直时上一章题解给 $a\propto t^{2/5}$, 所以 $T^2/\rho\propto a^3\propto t^{6/5}$. 原式在粒子熵自由度改变时 $aT$ 不恒定, 不能继续把 $A$ 视作同一个常数.
\end{solution}
